% !TEX root = cosine_is_human.tex
\appendix
\clearpage

\noindent{\Large\bfseries A\quad Model naming and factor values}
\applabel{A}{app:naming}
\vspace{1em}

Throughout the paper we identify each evaluated model by a four-part name of the form \texttt{\{Backbone\}-\{Adaptation\}-\{Head\}-\{Dataset\}}. Table~\ref{tab:naming} enumerates the possible values each component can take. Backbones are the feature extractors. Adaptation describes whether the backbone is frozen, partially trained via LoRA adapters, fully retrained, or used with externally-supplied pretrained weights. The head is the module that converts two (or three) backbone outputs into a similarity score; nine variants are compared including plain cosine distance, three cognitive-science-motivated architectures (Cog-v1 through Cog), the published Tversky Neural Network of \citet{doumbouya2025tversky}, the DreamSim margin-hinge ensemble loss, and a VLM logit-ratio judge. Dataset values record which NIGHTS subset a model was trained and evaluated on; omitting the dataset suffix defaults to NIGHTS-100k, our primary evaluation target.

\begin{table}[!ht]
\caption{Model naming scheme. Every model name in the paper is a concatenation of one value from each of the four rows. Example: \texttt{DINOv3-LoRA-Cog} uses DINOv3 as backbone, rank-8 LoRA on qkv/proj, our Cog head (frozen cosine alignment + GCM context + Tversky contrast), evaluated on the full NIGHTS-100k corpus.}
\label{tab:naming}
\centering
\small
\begin{tabular}{lll}
\toprule
Component & Values & Description \\
\midrule
Backbone & DINOv1, DINOv2, DINOv3 & Self-supervised ViTs \\
 & C-RADIOv4 & Distillation of DINOv3+SigLIP2+SAM3 \\
 & Qwen3-VL-4B & Vision-language model \\
 & DreamSim & DINO+CLIP+OpenCLIP ensemble \\
 & AlexNet-lin & LPIPS classical backbone \\
\midrule
Adaptation & ZS & Zero-shot, frozen backbone \\
 & LoRA & Rank-8 LoRA on qkv/proj + head trained \\
 & Pretrained & External weights (trained on NIGHTS-20k) \\
 & Scratch & Retrained on the stated dataset from scratch \\
\midrule
Head & CLS-cos & Cosine distance on CLS token \\
 & Tversky-frozen & Tversky contrast, $\alpha=\beta=0.5$ \\
 & Cog-v1 & Learned cross-attention + context + Tversky \\
 & Cog-v2 & Shared alignment + bottleneck + Tversky \\
 & Cog & Frozen cosine align + GCM context + Tversky \\
 & Cog-strat & Cog with inverse-frequency sampling \\
 & TNN & Doumbouya 2025 differentiable TNN \\
 & Hinge-ens & Margin hinge loss over ensemble features \\
 & VLM-judge & Qwen3-VL logit ratio on A/B choice \\
\midrule
Dataset & NIGHTS-100k & Full 100k, 64{,}420 non-tie, 35{,}580 ties held out \\
 & NIGHTS-100k-unan & X-0 subset, $N{=}20{,}017$ \\
 & NIGHTS-100k-weak & 2-1/3-1/4-1 subset, $N{=}36{,}476$ \\
 & NIGHTS-20k (published) & Filtered unanimous release \\
\bottomrule
\end{tabular}
\end{table}

\clearpage
\noindent{\Large\bfseries B\quad Per-fold raw numbers}
\applabel{B}{app:perfold}
\vspace{1em}

Table~\ref{tab:main} reports mean accuracy and standard deviation over five cross-validation folds. Table~\ref{tab:perfold} gives the accuracy of each of the twelve models on each fold. We form folds from the 64{,}420 non-tie triplets using fixed seed 42 and reuse the same indices for every model, so comparisons are not affected by data shuffling. For trained models, no fold differs from its model's mean by more than 1 percentage point, and the model rankings stay the same on every fold. Figure~\ref{fig:drop} plots each model's accuracy on the unanimous X-0 subset against its accuracy on the full corpus: every model drops 10.6--17.2 pp as the discarded disagreement is reincorporated.

\begin{figure}[ht]
\centering
\includegraphics[width=0.78\linewidth]{figures/drop_chart.pdf}
\caption{The same models use 5-fold cross-validation on unanimous X-0 test folds (blue; disjoint from training, comparable to published 20k) and all non-tie NIGHTS-100k triplets (red). Each drops 10 to 17 points because test sets differ. Dashed line: DreamSim's published 96.2\% on their 20k held-out split.}
\label{fig:drop}
\end{figure}

Figure~\ref{fig:agreement} gives the per-model curves behind Figures~\ref{fig:tier_ceiling} and~\ref{fig:gap_ties}, and Table~\ref{tab:tiergap} the accuracies by vote pattern.

\begin{figure}[!ht]
\centering
\includegraphics[width=0.78\linewidth]{figures/agreement_curve.pdf}
\caption{Accuracy versus human-agreement strength on NIGHTS-100k non-tie triplets. Accuracy rises with agreement for every model; DreamSim leads at most levels, but its released weights overlap our test folds. Unanimous patterns are pooled at 1.0; patterns with fewer than 100 triplets are omitted. Shaded band: the unanimous NIGHTS-20k subset.}
\label{fig:agreement}
\end{figure}

\begin{table}[!ht]
\caption{Accuracy by agreement tier, pooled over the five test folds, against the per-tier reproducibility ceiling (\S\ref{sec:ceiling}); the numbers behind Figures~\ref{fig:tier_ceiling} and~\ref{fig:gap_ties}. The LoRA columns are DINOv1-LoRA-Cog, DINOv2-LoRA-Cos, and DINOv3-LoRA-Cog. $^\dagger$Released weights overlap our test folds (Table~\ref{tab:main}).}
\label{tab:tiergap}
\centering
\footnotesize
\setlength{\tabcolsep}{4pt}
\begin{tabular}{lcc ccc ccc cc}
\toprule
 & & & \multicolumn{3}{c}{Zero-shot cosine} & \multicolumn{3}{c}{+ LoRA} & \multicolumn{2}{c}{DreamSim} \\
\cmidrule(lr){4-6}\cmidrule(lr){7-9}\cmidrule(lr){10-11}
Agreement & N & Ceiling & v1 & v2 & v3 & v1 & v2 & v3 & Pre$^\dagger$ & Scratch \\
\midrule
2-1 & 18{,}454 & 74\% & 62.3 & 59.4 & 60.5 & 61.7 & 63.9 & 64.1 & 65.2 & 64.7 \\
3-1 & 11{,}181 & 84\% & 67.6 & 64.9 & 66.5 & 68.3 & 72.1 & 72.0 & 73.0 & 71.9 \\
4-1 & 6{,}841 & 94\% & 71.5 & 68.3 & 69.8 & 72.8 & 77.3 & 77.7 & 79.5 & 77.2 \\
X-1 ($X \geq 5$) & 7{,}922 & 96.5--99.5\% & 76.2 & 72.7 & 73.8 & 78.8 & 82.7 & 83.0 & 85.2 & 83.5 \\
X-0 (unanimous) & 20{,}019 & 100\% & 88.6 & 83.9 & 85.7 & 90.7 & 93.8 & 94.5 & 97.9 & 94.1 \\
\midrule
All non-tie & 64{,}420 & $\approx$89\% & 74.1 & 70.6 & 72.0 & 75.1 & 78.3 & 78.7 & 80.7 & 78.7 \\
\bottomrule
\end{tabular}
\end{table}

\paragraph{Training-tier ladder.} Table~\ref{tab:ladder} asks whether the lower-agreement triplets help as training data. We trained DINOv2-LoRA-Cos on nested subsets of the fold-0 training split, ordered by agreement, and scored every run on the same fold-0 test split. A common validation set (10\% of the fold-0 training split) was removed from every run, and each run kept the epoch with the best validation accuracy on the tiers it trained on. Training on unanimous triplets alone gives 94\% of the gain over zero-shot that training on all triplets gives. Adding the remaining 21{,}380 triplets, almost all 2-1 and 3-1, to the first three tiers does not raise accuracy (78.4\% to 78.1\%). At equal size, unanimous triplets beat a random sample by 0.8 pp. In every run, validation accuracy peaked after one epoch, and by epoch 10 training accuracy had reached 100\% and test accuracy had fallen below zero-shot. These are single runs on one fold: an accuracy in the table has a standard error of 0.4 pp overall and 0.7--1.2 pp within a tier.

\begin{table}[!ht]
\caption{Training-tier ladder on fold 0. DINOv2-LoRA-Cos is trained on nested subsets of the training split and scored on the fold-0 test split (12{,}884 triplets) by vote pattern. The first row is the same backbone before training; the last is a random sample of all non-tie triplets at the size of the unanimous set.}
\label{tab:ladder}
\centering
\small
\begin{tabular}{lrcccccc}
\toprule
Training set & Triplets & 2-1 & 3-1 & 4-1 & X-1 ($X \geq 5$) & X-0 & All \\
\midrule
None (zero-shot) & 0 & 59.9 & 64.5 & 68.2 & 72.4 & 84.5 & 70.8 \\
X-0 & 14{,}459 & 63.7 & 71.4 & 76.2 & 81.0 & 93.1 & 77.7 \\
X-0, X-1 ($X \geq 5$) & 20{,}101 & 64.8 & 72.2 & 77.3 & 80.8 & 93.4 & 78.3 \\
X-0, X-1 ($X \geq 5$), 4-1 & 25{,}003 & 64.5 & 71.4 & 77.6 & 82.2 & 93.7 & 78.4 \\
All non-tie & 46{,}383 & 62.9 & 71.9 & 78.9 & 82.4 & 93.5 & 78.1 \\
\midrule
Random, X-0 size & 14{,}459 & 63.4 & 71.1 & 74.6 & 81.2 & 91.6 & 76.9 \\
\bottomrule
\end{tabular}
\end{table}


\begin{table}[!ht]
\caption{Per-fold accuracies (\%) on the NIGHTS-100k non-tie corpus ($N{=}64{,}420$ triplets; 12{,}884 per test fold; 5-fold cross-validation, seed 42). Each row is one model identified by the naming scheme of Table~\ref{tab:naming}. Columns are the five test folds; the same fold definitions are reused across all models. Mean and standard deviation over the five folds are reported in Table~\ref{tab:main} of the main body.}
\label{tab:perfold}
\centering
\small
\begin{tabular}{lccccc}
\toprule
Model & Fold 1 & Fold 2 & Fold 3 & Fold 4 & Fold 5 \\
\midrule
LPIPS-AlexNet-Pretrained & 65.4 & 65.5 & 65.3 & 65.3 & 65.0 \\
DINOv1-ZS-CLS-cos & 74.8 & 74.0 & 73.8 & 73.6 & 74.2 \\
DINOv2-ZS-CLS-cos & 70.8 & 71.0 & 70.8 & 70.3 & 70.0 \\
DINOv3-ZS-CLS-cos & 72.4 & 72.3 & 72.1 & 71.4 & 71.7 \\
DINOv1-LoRA-Cog & 74.8 & 75.6 & 75.0 & 75.3 & 74.9 \\
DINOv2-LoRA-Cog & 77.4 & 77.5 & 76.2 & 77.4 & 77.1 \\
DINOv2-LoRA-Cos & 78.7 & 78.4 & 77.6 & 78.3 & 78.8 \\
DINOv2-LoRA-Cog-strat & 76.4 & 76.4 & 76.2 & 75.8 & 77.2 \\
DINOv2-LoRA-TNN & 77.2 & 77.5 & 76.7 & 77.4 & 78.1 \\
DINOv3-LoRA-Cog & 78.7 & 78.9 & 78.3 & 78.4 & 79.1 \\
DreamSim-Pretrained & 80.5 & 80.9 & 80.4 & 80.9 & 80.8 \\
DreamSim-Scratch & 79.4 & 79.7 & 77.9 & 78.4 & 78.2 \\
\bottomrule
\end{tabular}
\end{table}

\clearpage
\noindent{\Large\bfseries C\quad Full Cohen's $\kappa$ matrix}
\applabel{C}{app:kappa}
\vspace{1em}

Cohen's $\kappa$ quantifies the agreement between two classifiers after correcting for the agreement expected by chance. We compute $\kappa$ pairwise between eleven of the twelve models in Table~\ref{tab:main} (the TNN head is omitted) on the pooled 64{,}420 non-tie test predictions (every triplet appears in exactly one fold's test split, so the pooled set is a disjoint union of the five held-out folds). Values near 1 indicate redundant models that would gain little from ensembling; values near 0.4--0.6 indicate models that make different errors and are candidates for ensembling. In Figure~\ref{fig:kappa} the three DINO zero-shot entries cluster tightly ($\kappa>0.63$), as do the five LoRA-adapted DINO variants ($\kappa>0.6$), confirming that the DINO family makes highly correlated errors. LPIPS sits further out ($\kappa=0.25\text{--}0.36$ with any DINO-family model), which is consistent with its different underlying feature basis. DreamSim-Pretrained and DreamSim-Scratch reach $\kappa=0.75$, the highest pairwise agreement in the matrix: retraining the DreamSim ensemble on NIGHTS-100k does not shift its error structure, it only shifts the overall accuracy by 2 pp. The scratch model is structurally the same predictor, trained on a cleaner subset, that fails on the same ambiguous triplets.

\begin{figure}[!ht]
\centering
\includegraphics[width=0.55\linewidth]{figures/kappa_matrix.pdf}
\caption{Cohen's $\kappa$ heatmap across eleven of the twelve models (TNN omitted) on NIGHTS-100k non-tie predictions ($N=64{,}420$). Each cell is the $\kappa$ between the row and column model's binary left/right predictions on the same set of triplets. Diagonal is 1 by construction. $\kappa=1$ is perfect agreement; $\kappa=0$ is chance. Two blocks stand out: the DINO-LoRA block (within-block $\kappa>0.6$), and the two DreamSim variants that agree with each other at $\kappa=0.75$ but with LPIPS at only 0.25--0.36.}
\label{fig:kappa}
\end{figure}

\paragraph{Is $\kappa$ sufficient when labels have limited reproducibility?} One reviewer concern is that pooled $\kappa$ mixes agreement on reproducible-label triplets (X-0, ceiling 100\%) with agreement on low-reproducibility triplets (2-1, ceiling 74\%). If two models ``agree'' on a 2-1 triplet, are they agreeing on shared signal or shared noise? We address this by stratifying $\kappa$ by agreement pattern: computing $\kappa$ separately on the X-0 (unanimous), X-1 (strong, $X\geq 5$), and weak (2-1/3-1/4-1) subsets. Table~\ref{tab:kappa-strat} shows that the ranking of redundancy is preserved across all strata, with magnitudes declining in the expected way on weak-majority triplets where label noise is higher.

\begin{table}[!ht]
\caption{Cohen's $\kappa$ stratified by human-agreement pattern on the pooled NIGHTS-100k non-tie predictions. The ranking of redundancy (DreamSim-Pretrained $\leftrightarrow$ DreamSim-Scratch highest, LPIPS $\leftrightarrow$ DreamSim lowest) is identical at every stratum. On the unanimous subset, where labels have 100\% reproducibility ceiling, the DreamSim-Pretrained $\leftrightarrow$ DreamSim-Scratch agreement rises to $\kappa=0.90$, the strongest piece of evidence that retraining the ensemble on NIGHTS-100k does not alter its decision surface.}
\label{tab:kappa-strat}
\centering
\small
\resizebox{\linewidth}{!}{%
\begin{tabular}{lcccc}
\toprule
Pair & Unanimous X-0 & Strong X-1 & Weak (2-1/3-1/4-1) & Pooled \\
     & $(N=20{,}005)$ & $(N=7{,}922)$ & $(N=36{,}476)$ & $(N=64{,}420)$ \\
\midrule
DreamSim-Pretrained $\leftrightarrow$ DreamSim-Scratch & 0.90 & 0.76 & 0.68 & 0.75 \\
DreamSim-Pretrained $\leftrightarrow$ DINOv3-LoRA-Cog & 0.89 & 0.69 & 0.58 & 0.69 \\
DINOv2-LoRA-Cos $\leftrightarrow$ DINOv3-LoRA-Cog    & 0.85 & 0.64 & 0.55 & 0.67 \\
DINOv3-ZS $\leftrightarrow$ DreamSim-Pretrained         & 0.72 & 0.52 & 0.44 & 0.54 \\
LPIPS $\leftrightarrow$ DreamSim-Pretrained              & 0.52 & 0.34 & 0.27 & 0.30 \\
\bottomrule
\end{tabular}}
\end{table}

The pooled $\kappa$ in Figure~\ref{fig:kappa} and in the main body is therefore a weighted average over agreement patterns with different label reliability. The stratified view shows two things: first, the redundancy claim (DreamSim variants and DINO-LoRA pairs are structurally similar classifiers, LPIPS is independent) holds on the subset where labels are authoritative; second, the lower $\kappa$ magnitudes on weak-majority patterns partly reflect the reduced maximum $\kappa$ achievable when classifier outputs are compared against noisy labels. Both observations are consistent with the \S\ref{sec:ceiling} reproducibility-ceiling analysis: on weak-majority triplets there is less reproducible signal for any classifier (or pair of classifiers) to latch onto.

\clearpage
\noindent{\Large\bfseries D\quad Additional experiments}
\applabel{D}{app:extra}
\vspace{1em}

This section documents experiments that did not make Table~\ref{tab:main} but took substantial compute: alternative similarity heads on frozen DINOv3 CLS embeddings, a multi-teacher distillation backbone (C-RADIOv4), a vision-language-model-as-judge (Qwen3-VL), ClickMe attention harmonization, and qualitative traces that make the agreement-stratified story visible on individual triplets. The BAPPS ceiling is reported in the main body (\S\ref{sec:bapps}).

\subsection*{D.0\quad Harmonized DINOv3 (attention alignment on ClickMe)}

\citet{fel2022harmonization} propose training a vision backbone to match human attention maps collected via ClickMe~\citep{linsley2019clickme}. The hypothesis is that attention-aligned features yield similarity-aligned embeddings. We tested this on DINOv3 with the authors' pretrained ClickMe attention predictor: 5-fold CV on NIGHTS-100k non-tie (matched splits to Table~\ref{tab:main}) gives Harmonized-DINOv3 $= 64.0 \pm 0.2\%$, an 8 pp drop from DINOv3 zero-shot on the same folds (Table~\ref{tab:harm}). The near-zero across-fold variance indicates the drop is not noise.

The interpretation in the Discussion is the level-of-abstraction mismatch: ClickMe captures where humans look during categorization, not which compositional features drive perceived ensemble-image similarity. The result is consistent with attention-alignment being a useful signal for LPIPS-style pixel-distortion metrics but not for the mid/high-level discrimination NIGHTS requires.

\begin{table}[ht]
\caption{Harmonized DINOv3 on NIGHTS-100k, 5-fold CV. Zero-shot DINOv3 reproduced from Table~\ref{tab:main} for direct comparison on the same test splits.}
\label{tab:harm}
\centering
\small
\begin{tabular}{lccccccc}
\toprule
Model & F1 & F2 & F3 & F4 & F5 & Mean $\pm$ Std & $\Delta$ \\
\midrule
DINOv3-ZS-CLS-cos            & 72.4 & 72.3 & 72.1 & 71.4 & 71.7 & $72.0 \pm 0.4$ & --     \\
Harmonized DINOv3            & 63.8 & 64.1 & 63.9 & 64.4 & 63.9 & $64.0 \pm 0.2$ & $-8.0$ \\
\bottomrule
\end{tabular}
\end{table}

\subsection*{D.1\quad Alternative similarity heads on frozen CLS}

The main body establishes that cosine distance on DINO CLS tokens is the strongest zero-shot signal on NIGHTS-100k. A natural follow-up question is whether any other similarity function on the same frozen embedding improves that number. We swept ten alternatives spanning basic distances, kernel methods, non-Euclidean geometries, distribution-based interpretations, and learned metric-learning heads, all evaluated on the same 5-fold cross-validation on NIGHTS-100k non-tie (Table~\ref{tab:head-sweep}).

\begin{table}[!ht]
\caption{Alternative similarity heads on frozen DINOv3 CLS embeddings, 5-fold CV on NIGHTS-100k non-tie ($N=64{,}420$). Cosine is the baseline from Table~\ref{tab:main}. ``Learned'' heads are trained with pairwise BCE for 10 epochs; all other heads are parameter-free. No alternative beats plain cosine by more than 0.1 pp, within fold-to-fold s.d.; learned heads with more parameters do systematically worse. Retrieval-style losses did not help either (cosine-softmax-ranking and margin-triplet variants of cosine all remained within 0.1 pp of plain cosine on a fold-0 pilot).}
\label{tab:head-sweep}
\centering
\small
\begin{tabular}{llrr}
\toprule
Head & Type & Trainable params & Accuracy \\
\midrule
\textbf{Cosine} (baseline)                & parameter-free, angular & 0        & \textbf{72.01 $\pm$ 0.4} \\
Neg-dot product                           & parameter-free, inner   & 0        & 72.01 $\pm$ 0.3 \\
Arc-cosine kernel (order-1)               & Cho \& Saul 2009        & 0        & 72.11 $\pm$ 0.3 \\
Poincar\'e hyperbolic distance            & non-Euclidean           & 0        & 72.09 $\pm$ 0.3 \\
L2 Euclidean                              & parameter-free, metric  & 0        & 71.56 $\pm$ 0.3 \\
L1 Manhattan                              & parameter-free, metric  & 0        & 71.50 $\pm$ 0.2 \\
JSD on softmax(CLS)                       & distribution divergence & 0        & 71.37 $\pm$ 0.2 \\
\midrule
Mahalanobis diag (learned)                & learned metric          & 768      & 71.79 $\pm$ 0.2 \\
Mahalanobis rank-32 (learned)             & learned metric          & 24{,}576 & 70.26 $\pm$ 0.4 \\
Learned bilinear rank-32                  & learned metric          & 49{,}152 & 68.35 $\pm$ 0.3 \\
Siamese MLP on $[a, b, |a{-}b|, a{\odot}b]$ & learned similarity    & $\sim$800k & 68.39 $\pm$ 0.5 \\
\bottomrule
\end{tabular}
\end{table}

Cosine, neg-dot, arc-cosine, and Poincar\'e differ by at most 0.1 pp, less than the fold-to-fold standard deviation. L2, L1, and JSD trail cosine by $\sim$0.5--0.7 pp. Adding parameters makes performance worse: a Mahalanobis metric with 24{,}576 extra parameters loses 1.7 pp, while an 800k-parameter Siamese MLP loses 3.6 pp despite access to both embeddings. The learned heads appear to fit fold-specific label noise (\S\ref{sec:ceiling}) rather than learn a better metric. On frozen DINO CLS embeddings, no tested head beats cosine by more than 0.1 pp, and all learned heads perform worse.

\subsection*{D.2\quad Multi-teacher distillation (negative result)}

C-RADIOv4-SO400M~\citep{conwell2024large} distills DINOv3, SigLIP2, and SAM3 into a single ViT. We hypothesized that a model with stronger semantic, linguistic, and segmentation priors would produce better perceptual similarity. We applied rank-8 LoRA to qkv/proj and a cosine head on the summary token and trained for one epoch on fold 0 (approximately 28 hours at the model's required 1008$\times$1008 input resolution for its RoPE embeddings).

Epoch-1 test accuracy on fold 0 was 58.2\%, below every DINO backbone in Table~\ref{tab:main} by 13 pp or more, and below LPIPS by 7 pp. We stopped after one epoch. The model's pretraining objectives mix spatial, semantic, and linguistic signals that appear to dilute the fine-grained synthetic-image similarity signal NIGHTS requires. Training at 224$\times$224 is not possible without retraining the RoPE embeddings, and the full-resolution compute cost makes multi-fold cross-validation impractical for this backbone at our compute budget.

\subsection*{D.3\quad VLM-as-judge (negative result)}

Modern vision-language models can describe image content at the scene level. We tested whether a VLM can directly answer the 2AFC question. Using Qwen3-VL-4B-Instruct on a 100-triplet sample from fold 0, we constructed a prompt showing the reference and two candidates with randomized A/B order, and extracted the logit ratio of the A vs.\ B tokens at the first generated position. We applied position randomization to remove A/B bias.

\begin{table}[!ht]
\caption{Qwen3-VL-4B-Instruct on a 100-triplet sample of NIGHTS-100k fold 0. The stratified rows in the submitted version were incorrect and are corrected here.}
\label{tab:qwen}
\centering
\small
\begin{tabular}{lcc}
\toprule
Triplet subset & N & Accuracy \\
\midrule
Unanimous (X-0)            & 26 & 69.2\% \\
Strong-majority (X-1, $X{\geq}5$) & 13 & 84.6\% \\
Weak-majority (2-1/3-1/4-1) & 61 & 55.7\% \\
\midrule
All                        & 100 & 63.0\% (Wilson 95\% CI 53.2--71.8) \\
\bottomrule
\end{tabular}
\end{table}

Qwen3-VL reaches 63.0\% overall, below LPIPS and every zero-shot backbone in Table~\ref{tab:backbones}. On unanimous triplets it reaches 69.2\%, while every zero-shot backbone exceeds 75\%, and it is near chance on weak-majority triplets. Each agreement stratum has only 13 to 61 triplets, so only the overall figure should be read quantitatively. The model can reason about scene content but does poorly on the fine-grained similarity judgments NIGHTS asks for.

\clearpage
\noindent{\Large\bfseries D.4\quad Model confidence versus the human vote split}
\applabel{D.4}{app:confidence}
\vspace{1em}

The headline result is in \S\ref{sec:confidence}. We map each model's score to a probability $P(\text{left})$ with out-of-fold Platt scaling (fit on four folds, applied to the fifth), define model confidence as $2|P-0.5|$, and compare it to the confidence the human vote supports, $2|v-0.5|$ for left-vote fraction $v$; the 35{,}580 tied triplets never enter the fit. Expected calibration error (ECE) is reported against both the binarized majority label and the human vote fraction. Figure~\ref{fig:confcal} reports the full picture: confidence by agreement pattern against the confidence the votes themselves support, a reliability curve of calibrated $P(\text{left})$ against the observed vote fraction, and mean confidence on the 35{,}580 split-vote triplets. Figure~\ref{fig:confcum} repeats the measurement cumulatively, starting from the published 20k and relaxing the agreement filter one vote pattern at a time out to the full 100k. The confidence the votes support falls from 1.000 to 0.414 across that sweep; DreamSim-Pretrained falls from 0.790 to 0.563, and the two curves cross near 70{,}000 triplets.

\begin{table}[ht]
\caption{Mean model confidence $2|P(\text{left})-0.5|$ by agreement pattern after out-of-fold Platt scaling, against the confidence the vote split supports. X-0 pools the 2-0 to 10-0 patterns (the one 1-0 and two 11-0 triplets are excluded).}
\label{tab:confidence}
\centering
\small
\begin{tabular}{lrcccc}
\toprule
Agreement & $N$ & Vote-supported & DreamSim-Pretrained & DINOv1-ZS & LPIPS \\
\midrule
50/50 (tie) & 35{,}580 & 0.000 & 0.484 & 0.383 & 0.251 \\
2-1         & 18{,}454 & 0.333 & 0.492 & 0.386 & 0.250 \\
4-1         & 6{,}841  & 0.600 & 0.544 & 0.412 & 0.260 \\
X-0         & 20{,}016 & 1.000 & 0.790 & 0.580 & 0.324 \\
\bottomrule
\end{tabular}
\end{table}

\begin{figure}[ht]
\centering
\includegraphics[width=\linewidth]{figures/confidence_calibration.pdf}
\caption{Model confidence against the human vote split on NIGHTS-100k. (A) confidence by agreement pattern; black is the confidence the observed vote split supports. (B) reliability against the observed vote fraction; legend gives ECE against the vote fraction and against the binarized majority label. (C) mean confidence on the 35{,}580 triplets where humans split 50/50. Gray bars are LoRA-trained heads reporting their native sigmoid.}
\label{fig:confcal}
\end{figure}

\begin{figure}[ht]
\centering
\includegraphics[width=\linewidth]{figures/confidence_by_agreement.pdf}
\caption{The same measurement on the agreement split used throughout the paper. (A) the 17 raw vote patterns on the majority-vote-fraction axis of Figure~\ref{fig:agreement}; the unanimous block is the published 20k. (B) cumulative corpus, relaxing the filter one pattern at a time from the published 20k to the full 100k.}
\label{fig:confcum}
\end{figure}

\clearpage
\noindent{\Large\bfseries D.5\quad Per-axis stimulus manipulation results}
\applabel{D.5}{app:axes_detail}
\vspace{1em}

These are the per-axis results of the Figure-6-style replica described in \S\ref{sec:shape_axes}. We report them here, and not in the main body, because NIGHTS provides no human judgments on manipulated stimuli, and we collected none: the orientation, color, shading and blur axes therefore measure model invariance only and cannot be set against human behavior. The foreground/background axis is treated separately in the main text, where IN-9 and ObjectNet supply an external anchor.

\begin{figure}[ht]
\centering
\includegraphics[width=0.88\linewidth]{figures/fig6_ablation_heatmap.pdf}
\caption{DreamSim-style Figure-6 ablation replica. Accuracy drop (pp) relative to clean for each backbone $\times$ ablation. Bluer is more invariant; redder is more fragile. DINOv3 zero-shot is the most invariant overall; DreamSim-Pretrained and our own DINOv2+LoRA are among the most fragile despite the best clean accuracy. 2000 unanimous NIGHTS triplets.}
\label{fig:fig6_replica}
\end{figure}

\begin{enumerate}[leftmargin=*,itemsep=1pt,topsep=1pt]
\item \textbf{DINOv3 was the most invariant zero-shot backbone} on every axis except background removal. Rotation drop was $\leq 0.7$ pp, grayscale drop was only 0.3 pp, mild blur actually \emph{helped} (+1.25 pp at $\sigma=4$), and the foreground-only drop was 1.1 pp. Its clean accuracy of 85.4\% was below DINOv1's 88.5\% and DreamSim-Pretrained's 93.8\%, but its accuracy changed less under the tested manipulations.
\item \textbf{DreamSim-Pretrained traded invariance for clean-image accuracy.} It had the best clean score in the table ($93.8\%$) but the worst or near-worst drop on rotation ($-4.3$), color ($-5.7$), heavy blur ($-9.6$), foreground ($-9.05$), and background ($-14.5$). Margin-hinge LoRA training narrows the feature geometry around the triplet distribution, which produces the best in-distribution score at the cost of cross-manipulation stability.
\end{enumerate}

Both observations are statements about models. Neither supports a claim about the difference between model error and human error on these stimuli, and the revision no longer makes one.

\begin{table}[!ht]
\caption{Zero-shot DINOv2 with cosine on NIGHTS-100k (5-fold mean $\pm$ s.d.). Accuracy falls monotonically with size. The 70.1\% reference is the split-half agreement of the NIGHTS votes, the probability that the two halves of a triplet's votes choose the same majority; it is a conservative companion to the $\approx$89\% fresh-panel ceiling of \S\ref{sec:ceiling}.}
\label{tab:scale}
\centering
\small
\begin{tabular}{lcccc}
\toprule
DINOv2 & Params & NIGHTS-100k & vs.\ 70.1\% & Unanimous X-0 \\
\midrule
ViT-S/14 & 21M  & $72.2 \pm 0.5$ & $+2.1$ & 85.8 \\
ViT-B/14 & 86M  & $70.6 \pm 0.4$ & $+0.5$ & 83.9 \\
ViT-L/14 & 300M & $69.0 \pm 0.2$ & $-1.1$ & 81.8 \\
ViT-g/14 & 1.1B & $68.4 \pm 0.3$ & $-1.7$ & 80.5 \\
\bottomrule
\end{tabular}
\end{table}

\clearpage
\noindent{\Large\bfseries D.6\quad Zero-shot backbone sweep}
\applabel{D.6}{app:backbones}
\vspace{1em}

Every backbone below is evaluated with plain cosine on the post-norm CLS (or pooled) embedding at 224$\times$224, with no training, on the same five folds as Table~\ref{tab:main}. The DINOv2 size sweep is repeated in Table~\ref{tab:scale}.

\begin{table}[ht]
\caption{Zero-shot cosine accuracy on NIGHTS-100k non-tie (5-fold mean $\pm$ s.d., $N=64{,}420$) and on the unanimous X-0 subset.}
\label{tab:backbones}
\centering
\small
\begin{tabular}{llcc}
\toprule
Backbone & Family & NIGHTS-100k & Unanimous X-0 \\
\midrule
LPIPS-AlexNet          & CNN, supervised       & $65.3 \pm 0.2$ & 75.9 \\
ResNet-152             & CNN, supervised       & $67.1 \pm 0.4$ & 78.9 \\
CLIP ViT-B/16          & ViT, language-image   & $68.0 \pm 0.3$ & 80.0 \\
ResNet-50              & CNN, supervised       & $69.7 \pm 0.3$ & 82.4 \\
BiT-M R101x3           & CNN, supervised       & $70.2 \pm 0.3$ & 83.5 \\
BagNet-33              & CNN, bio-inspired     & $70.3 \pm 0.5$ & 83.6 \\
ConvNeXt-Large         & CNN, supervised       & $71.4 \pm 0.3$ & 85.4 \\
CORnet-S               & CNN, bio-inspired     & $71.7 \pm 0.2$ & 85.7 \\
VOneResNet-50          & CNN, bio-inspired     & $71.8 \pm 0.6$ & 85.7 \\
Supervised ViT-B/16    & ViT, supervised       & $72.4 \pm 0.2$ & 86.5 \\
\midrule
DINOv2 ViT-g/14        & ViT, self-supervised  & $68.4 \pm 0.3$ & 80.5 \\
DINOv2 ViT-L/14        & ViT, self-supervised  & $69.0 \pm 0.2$ & 81.8 \\
DINOv2 ViT-B/14        & ViT, self-supervised  & $70.6 \pm 0.4$ & 83.9 \\
DINOv3 ViT-B/16        & ViT, self-supervised  & $72.0 \pm 0.4$ & 85.7 \\
DINOv2 ViT-S/14        & ViT, self-supervised  & $72.2 \pm 0.5$ & 85.8 \\
DINOv1 ViT-B/16        & ViT, self-supervised  & $74.1 \pm 0.5$ & 88.6 \\
\midrule
DreamSim ensemble, untrained & DINO+CLIP+OpenCLIP & $74.5 \pm 0.4$ & 89.2 \\
\bottomrule
\end{tabular}
\end{table}

\noindent{\Large\bfseries D.7\quad Color control for the fine-tuning gain}
\applabel{D.7}{app:colorcontrol}
\vspace{1em}

To test whether the 4.2 pp that training adds over the untrained DreamSim ensemble (\S\ref{sec:gain}) comes from low-level color, we computed for every one of the 64{,}420 non-tie triplets the $\chi^2$ distance between the RGB histograms of each candidate and the reference, and split the triplets by whether that gap favors the human-majority candidate. The training gain was $+3.9$ pp where color could not discriminate and $+4.9$ pp where it could (95\% CI $\pm 0.7$ pp), so the gain is nearly flat across the split. A logistic regression on twelve pixel statistics alone reached 62.3\%, below every zero-shot backbone in Table~\ref{tab:backbones}.

\clearpage
\noindent{\Large\bfseries E\quad Training details}
\applabel{E}{app:training}
\vspace{1em}

This section documents the full training recipe for reproducibility. All trained models in Table~\ref{tab:main} use the same pipeline except where explicitly noted; DreamSim-Scratch uses the authors' published training code verbatim and is documented separately at the end.

\paragraph{Preprocessing.} Images are decoded from the NIGHTS release JPEGs, resized to $224\times 224$ with bicubic interpolation, and center-cropped. Normalization follows the backbone: ImageNet statistics (mean $[0.485, 0.456, 0.406]$, std $[0.229, 0.224, 0.225]$) for DINOv1, DINOv2, and DINOv3; $0.5/0.5$ for C-RADIOv4; the processor-supplied normalization for DreamSim and Qwen3-VL. No data augmentation is used at training time: the triplet task is orientation-sensitive, and the reference/candidate pairing is the supervision signal.

\paragraph{LoRA configuration.} LoRA adapters have rank 8, $\alpha = 16$, and dropout 0.05, and are inserted on the \texttt{qkv} and \texttt{proj} linear layers of every transformer block. The backbone's final LayerNorm is also unfrozen. For DINOv2 ViT-B/14 with 12 blocks this yields 442{,}368 LoRA parameters plus 1{,}536 LayerNorm parameters, which is approximately 0.5\% of the 86M backbone. The DreamSim ensemble retraining uses rank 16 on the same modules to match the authors' released configuration.

\paragraph{Heads.} Zero-shot models use cosine distance on the post-norm CLS token at $224\times224$ input. The Cog head has 197k parameters: frozen cosine alignment, GCM context, and Tversky contrast with learnable $\alpha, \beta$. The Tversky Neural Network of \citet{doumbouya2025tversky} is matched at 197k parameters, with product intersection and subtract-match difference functions and 256 features.

\paragraph{Optimizer.} Schedule-Free AdamW \citep{defazio2024schedulefree} with $\beta_1 = 0.9$, $\beta_2 = 0.999$, weight decay $10^{-2}$, 100 linear warm-up steps, and a constant learning rate thereafter (no cosine schedule). Separate learning rates are used: $10^{-4}$ for LoRA and layer-norm parameters, $2\times 10^{-4}$ for the head. Gradients are clipped to a global norm of 1.0 before each optimizer step.

\paragraph{Batch and compute.} Batch size is 32 triplets per step (each triplet contributes three images to the forward pass). Training runs for 10 epochs over the fold's training split with batch size 32. Compute is a shared cluster of NVIDIA RTX A6000 GPUs (48\,GB). A full 5-fold run of DINOv2-LoRA-Cog takes approximately 3--5 hours when five GPUs run concurrently, one per fold.

\paragraph{Checkpoint selection.} A held-out validation split of 10\% of each fold's training data selects the best-of-10-epochs checkpoint by validation accuracy; the selected checkpoint is evaluated once on the fold's test split, and that single number appears in Table~\ref{tab:perfold}. Per-triplet test-set predictions are persisted for every (model, fold) pair so that the Cohen's $\kappa$ matrix, McNemar tests, and error-distribution analyses can be reproduced without re-running training.

\paragraph{DreamSim retraining from scratch.} We run the authors' published training code unchanged with the released ensemble configuration: DINO+CLIP+OpenCLIP, rank-16 LoRA, margin-0.05 hinge loss, batch size 16, 6 epochs. The release filters training triplets to those with at least 6 votes, so the scratch model sees only strong-majority triplets during training; we evaluate on the same full non-tie test fold as every other model in Table~\ref{tab:main}. This mismatch is the point of the comparison, not an oversight: it isolates what the authors' training recipe would produce when the 100k distribution is available.

\clearpage
\noindent{\Large\bfseries F\quad Shape-only class separability and retrieval}
\applabel{F}{app:shape_only}
\vspace{1em}

\begin{table}[ht]
\caption{Shape-only class separability. Leave-one-out nearest-prototype accuracy on the Geirhos filled-silhouette and edge stimuli. Chance is 6.25\% (16 classes). Referenced from \S\ref{sec:shape_axes}.}
\label{tab:shape_only}
\centering
\small
\begin{tabular}{lcccc}
\toprule
Backbone & Silh 1-NN & Silh proto & Edges 1-NN & Edges proto \\
\midrule
ResNet-50        & 77.5 & 81.3 & 63.1 & 67.5 \\
DINOv1           & 71.3 & 78.1 & 66.3 & 68.1 \\
CLIP-ViT-B/16    & 73.1 & 78.8 & 79.4 & 81.3 \\
DINOv2           & 81.3 & 86.3 & 71.3 & 80.0 \\
DINOv3           & 82.5 & 87.5 & 86.9 & 91.3 \\
ConvNeXt-Large   & 84.4 & 87.5 & 90.0 & 93.8 \\
\bottomrule
\end{tabular}
\end{table}

\begin{figure}[ht]
\centering
\includegraphics[width=\linewidth]{figures/silhouette_retrieval.png}
\caption{Shape-only nearest-neighbor retrieval. For four query silhouettes (bicycle, cat, elephant, keyboard), each backbone's top-3 nearest silhouettes by cosine distance in CLS space. Green outline: match; red: mismatch. Every backbone, including ResNet-50 with its near-zero cue-conflict shape-bias, retrieves predominantly within-class; shape representation is present in every backbone we tested.}
\label{fig:silh_retrieval}
\end{figure}

\begin{figure}[ht]
\centering
\includegraphics[width=0.8\linewidth]{figures/cost_of_accuracy.pdf}
\caption{The cost of clean-image accuracy. Each point is a backbone; vertical axis is the mean magnitude of accuracy drop across the 11 Fig-6 ablation conditions. Zero-shot backbones (circles) cluster in the low-fragility region; LoRA-trained models and DreamSim-Pretrained (diamonds) shift up and right. Dashed line is a least-squares trend.}
\label{fig:cost_acc}
\end{figure}

\paragraph{Tie-based bias audit.} On the 7,116-triplet tie subset held out per fold, LoRA-trained models predict left with frequency 49.5--51.1\% (no positional bias). Zero-shot DINOs show 49.6--50.8\% and DreamSim 51\%, also near chance. The one model with tie bias outside that range is Qwen3-VL-4B (55\% left, confidence 0.30), consistent with known VLM positional preferences. LoRA adaptation on NIGHTS-100k does not induce spurious tie preferences.

\clearpage
\noindent{\Large\bfseries G\quad Full 2AFC-vs-JND axis table}
\applabel{G}{app:axes}
\vspace{1em}

\begin{table}[ht]
\caption{Per-backbone 2AFC accuracy (Fig-6 clean, 2000 unanimous triplets) and JND AUC (1{,}216 reference--candidate pairs from 608 triplets, same/different). Extended form of Figure~\ref{fig:tasks_decouple}. Among the zero-shot backbones the two rankings are only weakly related (Spearman $\rho = 0.37$, $p = 0.19$); the two models trained on NIGHTS 2AFC lead on both. A single AUC carries a 95\% bootstrap interval of about $\pm 0.04$.}
\label{tab:axes}
\centering
\small
\begin{tabular}{lccl}
\toprule
Backbone & 2AFC acc (\%) & JND AUC & Family \\
\midrule
DINOv1              & 88.45 & 0.674 & DINO self-distill \\
DINOv2              & 84.65 & 0.677 & DINO self-distill \\
DINOv3              & 85.35 & 0.697 & DINO self-distill \\
CLIP-ViT-B/16       & 80.60 & 0.661 & image-text \\
OpenCLIP-ViT-B/16   & 83.95 & 0.696 & image-text \\
ViT-B/16 & 86.90 & 0.670 & label-supervised ViT \\
ConvNeXt-Large      & 85.40 & 0.648 & modern CNN \\
ResNet-50           & 84.05 & 0.647 & classical CNN \\
ResNet-152          & 83.45 & 0.653 & classical CNN \\
BiT-M R101x3        & 83.50 & 0.634 & JFT-scale CNN \\
CORnet-S            & 86.60 & 0.675 & biologically-inspired \\
BagNet-33           & 84.20 & 0.651 & biologically-inspired \\
VOneResNet50        & 85.55 & 0.674 & biologically-inspired \\
DreamSim-ZS     & 89.00 & 0.676 & concatenated ensemble, ZS \\
DINOv2 + LoRA       & 89.30 & 0.720 & LoRA-trained on NIGHTS \\
DreamSim-Pretrained & 93.75 & 0.727 & LoRA-trained ensemble \\
\bottomrule
\end{tabular}
\end{table}

\clearpage
\noindent{\Large\bfseries H\quad BAPPS-vs-NIGHTS quadrant view}
\applabel{H}{app:crossover}
\vspace{1em}

\begin{figure}[ht]
\centering
\includegraphics[width=0.45\linewidth]{figures/bapps_vs_nights_crossover.pdf}
\caption{Quadrant view of Figure~\ref{fig:crossover}. Dashed lines mark the 69.0\% (BAPPS) and 70.1\% (NIGHTS) reference levels of Table~\ref{tab:bapps}. The top-right quadrant lies above both. Two models fall inside it: DreamSim-Pretrained, after supervised training on the NIGHTS-20k subset, and zero-shot DINOv1, with no training.}
\label{fig:crossover_quadrant}
\end{figure}

\begin{figure}[ht]
\centering
\includegraphics[width=0.72\linewidth]{figures/bapps_vs_nights_spectrum.pdf}
\caption{Perceptual similarity is a spectrum. Each model's accuracy at low-level BAPPS (left) connects to its accuracy at mid-to-high-level NIGHTS-100k (right). Crossing lines show rank inversion across the spectrum: LPIPS is a BAPPS specialist that trails on NIGHTS; DINOv2/v3 invert the pattern. Two models score above both reference levels: DreamSim after supervised NIGHTS-20k training, and zero-shot DINOv1 with no training.}
\label{fig:crossover}
\end{figure}

\begin{table}[ht]
\caption{Split-half agreement for BAPPS 2AFC validation (5 votes per triplet): the probability that the two halves of a triplet's votes choose the same majority, a conservative companion to the fresh-panel repeatability of \S\ref{sec:ceiling} (68.3\% for 3-2, 94.2\% for 4-1). The lower rows give per-category 2AFC accuracy against the binary majority label; their unweighted mean is the BAPPS column of Table~\ref{tab:bapps}.}
\label{tab:bapps-ceiling}
\centering
\small
\begin{tabular}{lcccccc}
\toprule
Category & CNN & Traditional & Super-res & Color & Frame-int & Deblur \\
\midrule
N        & 4{,}720 & 4{,}720 & 10{,}856 & 4{,}720 & 1{,}888 & 9{,}440 \\
Ceiling  & 83.2\% & 79.1\% & 70.0\% & 64.1\% & 62.8\% & 59.6\% \\
\midrule
LPIPS-AlexNet           & 86.7 & 77.7 & 76.2 & 66.3 & 66.9 & 63.8 \\
DINOv1-ZS               & 85.2 & 80.0 & 73.2 & 62.0 & 66.1 & 62.4 \\
DINOv2-ZS               & 81.0 & 75.4 & 63.7 & 58.1 & 66.6 & 59.1 \\
DINOv3-ZS               & 83.2 & 77.2 & 65.8 & 59.8 & 64.4 & 60.2 \\
DreamSim-Pretrained & 86.8 & 82.0 & 75.0 & 67.7 & 67.5 & 62.6 \\
\bottomrule
\end{tabular}\\[3pt]
\textbf{All BAPPS: N=36{,}344, split-half agreement = 69.0\%.}
\end{table}

\begin{table}[ht]
\caption{Accuracy on ImageNet-prompted and non-ImageNet-prompted subsets of NIGHTS-100k non-tie. $\Delta$ is ImageNet minus non-ImageNet.}
\label{tab:imnet}
\centering
\small
\begin{tabular}{lccc}
\toprule
Model & ImageNet & Non-ImageNet & $\Delta$ \\
\midrule
LPIPS-AlexNet-Pretrained       & 64.6 & 66.1 & $-1.5$ \\
DINOv1-ZS-CLS-cos              & 74.5 & 73.6 & $+0.9$ \\
DINOv3-ZS-CLS-cos              & 73.0 & 70.7 & $+2.3$ \\
DINOv2-LoRA-Cos                & 78.7 & 77.9 & $+0.7$ \\
DINOv3-LoRA-Cog             & 79.4 & 77.9 & $+1.5$ \\
DreamSim-Pretrained        & 81.0 & 80.3 & $+0.8$ \\
DreamSim-Scratch           & 79.0 & 78.3 & $+0.7$ \\
\bottomrule
\end{tabular}
\end{table}

\begin{figure}[ht]
\centering
\includegraphics[width=0.82\linewidth]{figures/error_distribution.pdf}
\caption{Left: share of each model's NIGHTS-100k errors, stratified by the human-agreement pattern of the underlying triplet (weak 2-1/3-1/4-1, strong X-1 with $X\geq 5$, unanimous X-0). Right: fraction of each model's wrong predictions where at least one human voter agreed with the model. The best models are almost never alone against the humans (DreamSim 96.6\%, DINOv3-LoRA 92.0\%); their errors live where humans themselves were split.}
\label{fig:errdist}
\end{figure}



\clearpage
\noindent{\Large\bfseries I\quad Extended \S\ref{sec:shape_axes} figures: shape bias, hGRU collapse, LoRA trajectory, 2AFC vs.\ JND}
\applabel{I}{app:shape_axes_extra}
\vspace{1em}

These five figures were referenced from \S\ref{sec:shape_axes} but deferred here for length.

\begin{figure}[ht]
\centering
\includegraphics[width=0.68\linewidth]{figures/shape_bias_cue_conflict.pdf}
\caption{Cue-conflict shape bias on the Geirhos 2018 stimuli. DINOv1 preferences are texture-dominated (10.3\%); DINOv2 (23.2\%) and DINOv3 (37.3\%) are more than three times more shape-biased. ResNet-50 is the most texture-dominant (3.75\%); ConvNeXt-Large (19.75\%) and CLIP-ViT-B/16 (29.83\%) sit between. Zero-shot, 1200 cue-conflict images, 16 classes. Humans score $\sim$95\% on the same stimuli \citep{geirhos2018imagenettrained}.}
\label{fig:shape_bias_cc}
\end{figure}

\begin{figure}[ht]
\centering
\includegraphics[width=0.78\linewidth]{figures/hgru_spatial_sweep.pdf}
\caption{Radial-mask sensitivity sweep. Accuracy on 2000 unanimous triplets as a function of radial mask radius, keeping only the center disc (fg, left) or only the surround ring (bg, right). DINO backbones transition smoothly; ResNet-50+hGRU collapses to chance on any foreground radius in $[0.1, 0.5]$ and any background radius below $0.7$, revealing that the lateral-recurrence stage learns a globally coupled representation that fails under any spatial mask.}
\label{fig:hgru_sweep}
\end{figure}

\begin{figure}[ht]
\centering
% figures/attention_zs_vs_lora.pdf is a 10-triplet strip (445.3 pt per triplet); show toaster, hotdog and spoonbill, the three the caption discusses
\includegraphics[width=0.92\linewidth,height=0.47\textheight,keepaspectratio,trim=0 3566 0 0,clip]{figures/attention_zs_vs_lora.pdf}\\[2pt]
\includegraphics[width=0.92\linewidth,height=0.235\textheight,keepaspectratio,trim=0 0 0 4008,clip]{figures/attention_zs_vs_lora.pdf}
\caption{Three high-vote NIGHTS unanimous triplets where DINOv2 + LoRA flips the majority-choice decision. Green border = human-majority candidate; red border = LoRA's (wrong) choice. Each panel is overlaid with last-layer CLS attention and annotated with the $\cos(\mathrm{ref}, \cdot)$ values the models computed. The examples share two signatures. (i) LoRA's flips occur at $\cos$ margins as small as 0.005 (spoonbill: $0.905$ vs $0.910$), exactly the voter-noise-scale differences the reproducibility analysis of Section~\ref{sec:ceiling} warned about; (ii) attention geometry is nearly identical between ZS and LoRA, so the shortcut is encoded in the feature representation (sim(full, fg-only) drops consistently for LoRA) and not in where the model looks. Pairwise similarities computed with DINOv2 + rank-8 LoRA + cosine at 224$\times$224.}
\label{fig:attn_zs_vs_lora}
\end{figure}

\begin{figure}[ht]
\centering
\includegraphics[width=0.78\linewidth]{figures/lora_strength_sweep.pdf}
\caption{Controlled LoRA-strength sweep on a single backbone. DINOv2 + rank-8 LoRA trained for 0--8 epochs on NIGHTS-20k unanimous, each checkpoint evaluated on the 2000-triplet Fig-6 ablation suite. Gray dots: 16-backbone context from Figure~\ref{fig:cost_acc}. Phase 1 (epoch 0\,$\to$\,1) mirrors the cross-backbone invariance tax slope. Phase 2 (epochs 1--8) keeps climbing the fragility axis while clean accuracy reverts toward the zero-shot baseline. Within-backbone evidence that the invariance tax is causal, not a family-level confound.}
\label{fig:lora_strength}
\end{figure}

\begin{figure}[ht]
\centering
\includegraphics[width=0.78\linewidth]{figures/tasks_decouple_scatter_v2.pdf}
\caption{2AFC accuracy and JND AUC per backbone. Horizontal axis: NIGHTS 2AFC accuracy on the 2000 unanimous-triplet subset; vertical axis: JND AUC on the 1{,}216 reference--candidate pairs of the 608 JND triplets. The green arrows trace DINOv1 $\to$ v2 $\to$ v3: v1 leads on 2AFC by 3.1--3.8 pp, but the three lie within 0.03 AUC of each other on JND, inside the $\pm 0.04$ sampling interval. The two models trained on NIGHTS 2AFC (upper right) lead on both tasks.}
\label{fig:tasks_decouple}
\end{figure}


\clearpage
\noindent{\Large\bfseries J\quad Qualitative triplet traces}
\applabel{J}{app:traces}
\vspace{1em}

The analyses in the main body are aggregate. This section shows the individual stimuli behind them, drawn unmodified from the released corpus. Figure~\ref{fig:traces_agreement} illustrates what the agreement patterns of \S\ref{sec:ceiling} look like on single triplets, one per pattern from 2-1 to 7-0. A 2-1 triplet carries a label that a fresh panel of three annotators would overturn about a quarter of the time; a 7-0 triplet carries one that reproduces every time. Both contribute equally to a headline accuracy.

\begin{figure}[!ht]
\centering
\includegraphics[width=\linewidth]{figures/triplet_traces_by_agreement_v2.png}
\caption{One NIGHTS-2AFC triplet per vote pattern, from 2-1 to 7-0: the reference, the candidate the human majority chose (green), and the alternative (red), each with its vote count. Prompts and triplet ids are given above each triplet.}
\label{fig:traces_agreement}
\end{figure}

Figure~\ref{fig:traces_hard} shows the complementary case: unanimous triplets with nine or ten votes, where the label is as stable as NIGHTS labels get, and where DreamSim-Pretrained still selects the other candidate. It does so on 20 of the 3{,}126 such triplets. These are the errors the reproducibility ceiling does not explain, and they are 3.4\% of that model's disagreements with the majority label.

\begin{figure}[!ht]
\centering
\includegraphics[width=\linewidth]{figures/unanimous_hard_examples_v2.png}
\caption{Unanimous triplets with nine or ten votes that DreamSim-Pretrained, the strongest model in Table~\ref{tab:main}, gets wrong: it picks the candidate no voter chose (red). The eight with the most votes are shown; vote pattern, prompt, and triplet id are given above each triplet.}
\label{fig:traces_hard}
\end{figure}

Figure~\ref{fig:traces_jnd} shows stimuli carrying both NIGHTS labels. All 608 shared triplets come from the X-0 unanimous tier, so the 2AFC/JND comparison of \S\ref{sec:discussion} is a comparison on the easy subset and not the hard one. JND majorities judged both candidates different from the reference in 466 triplets (77\%). In 137 triplets exactly one candidate was judged the same, the chosen one in 103 and the rejected one in 34, and in the remaining 5 both were. Comparing the two JND vote fractions within a triplet, voters judged the chosen candidate less different than the rejected one on 42.3\% of triplets, could not separate the two on 39.6\%, and judged the chosen one more different on 18.1\%.

\begin{figure}[!ht]
\centering
\includegraphics[width=\linewidth]{figures/jnd_2afc_examples_v2.png}
\caption{Triplets carrying both human labels, two for each of the three cases. The label at the top of each candidate gives its 2AFC votes; the label at the bottom gives the JND majority verdict and its vote count.}
\label{fig:traces_jnd}
\end{figure}

\clearpage

\noindent{\Large\bfseries K\quad Where human voters disagree: splits within and conflicts across the two NIGHTS tasks}
\applabel{K}{app:disagree}
\vspace{1em}

NIGHTS carries two kinds of human disagreement, and a majority label shows neither. The first is within a task: voters split on the same item. The second is across tasks: the 2AFC label and the JND label of the same image pair point in opposite directions. This section reports both from the released votes.

\paragraph{Splits within 2AFC.} Of the 64{,}420 non-tie triplets, 36{,}476 (57\%) are weak majorities (2-1, 3-1 or 4-1) and 7{,}922 have a single dissent among five or more voters; a further 35{,}580 triplets are ties. Because voting stopped at the first dissent, a split triplet usually has only three or four voters, and a fresh panel of the same size reproduces a 2-1 majority only 74\% of the time (\S\ref{sec:ceiling}). The DreamSim paper reports the split only indirectly, as the roughly 20\% of triplets that survived its unanimity filter, and the public release contains only the unanimous triplets.

\paragraph{Splits within JND.} JND voters judged each reference--candidate pair ``same'' or ``different'': three voters for 1{,}184 of the 1{,}216 pairs and six for the other 32. Of the pairs, 643 (52.9\%) were unanimous and 573 (47.1\%) were split; none was tied. Applying the fresh-panel resampling of \S\ref{sec:ceiling} to the JND votes, a split JND label is reproduced 74.7\% of the time and the average JND label 88.1\%, the same reliability profile as 2AFC. The DreamSim paper takes the majority of the three votes as the label and reports no split rate. It describes keeping triplets in which humans found exactly one candidate identical; in the release that is 137 of the 608 triplets (23\%).

\paragraph{Conflicts across the two tasks.} Table~\ref{tab:disagree} cross-tabulates the two labels on the 1{,}216 pairs. A pair is consistent when the 2AFC winner is judged ``same'' or the loser ``different,'' and in conflict when the winner is judged ``different'' or the loser ``same.'' The conflicts sit almost entirely in one place: unanimous 2AFC winners that JND voters judged different from the reference. This holds whether the JND vote was split (249 winners) or unanimous (251 winners), so it is not produced by noise in three-vote JND labels. At the triplet level, JND voters judged neither candidate ``same'' in 466 of the 608 triplets (77\%): the unanimous 2AFC vote chose between two images that JND voters saw as different from the reference. Of the 137 triplets in which exactly one candidate was judged ``same,'' that candidate was the 2AFC loser in 34 (25\%). The share of winners judged different does not fall as 2AFC agreement rises: 81\% at 6-0 and 7-0, 82\% at 8-0, 84\% at 9-0, and 91\% at 10-0 (22 triplets). Figure~\ref{fig:traces_jnd} shows examples.

\begin{table}[ht]
\caption{The two human labels on the 1{,}216 reference--candidate pairs of the 608 NIGHTS-JND triplets, all unanimous in 2AFC. Columns give the JND vote, unanimous or split, grouped by its majority. Bold cells point against the 2AFC choice (chosen judged ``different,'' rejected judged ``same''); every candidate is a visibly distorted image, so the first kind is expected, and the informative comparison is between the two rows. The last column is each row's share pointing against the 2AFC choice.}
\label{tab:disagree}
\centering
\small
\begin{tabular}{lrrrrr}
\toprule
 & \multicolumn{2}{c}{JND majority ``different''} & \multicolumn{2}{c}{JND majority ``same''} & \\
\cmidrule(lr){2-3}\cmidrule(lr){4-5}
2AFC label & unanimous & split & split & unanimous & against 2AFC \\
\midrule
Winner (chosen), $n=608$ & \textbf{251} & \textbf{249} & 92 & 16 & 82.2\% \\
Loser (rejected), $n=608$ & 372 & 197 & \textbf{35} & \textbf{4} & 6.4\% \\
\midrule
All pairs, $n=1{,}216$ & 623 & 446 & 127 & 20 & 44.3\% \\
\bottomrule
\end{tabular}
\end{table}

\paragraph{What the conflict means.} 2AFC is a forced comparison with no ``neither'' option: a candidate wins by being closer to the reference than the alternative, however different both are. JND asks whether a single candidate can be told apart from the reference. A unanimous 2AFC label therefore certifies a ranking, not similarity, and the 77\% of triplets with both candidates judged different show that most of these unanimous labels rank two visibly different images. Scoring a model against the majority label of either task discards both the splits within each task and the conflict between them.

\clearpage
